% Template for Provadis University scientific works
% Based on "University of Ulm Internship Report Template"
% by Max Sch.
% CC BY 4.0
%
% Einstellungen.tex
%
% General settings are defined here.

%Packages are extensions of the actual program and facilitate workflows, representations, formatting etc.
%The following packages have been installed once for creating a bachelor thesis.
%They can be expanded/reduced as needed.

%%%%%%%%%%%%%%%%%%%%%%%%%%%%%%%%%%%%%%%%%%%%%%%%%%%%%%%%%
%%%%%%%%%%%%%%%%%%%%%%%%%%%%%%%%%%%%%%%%%%%%%%%%%%%%%%%%%

% Times New Roman Alternative, pdflatex compatible
\usepackage{newtxtext}

% Package for Times font
%\usepackage{mathptmx}

\usepackage[english]{babel}
% Package for English language (translations from Chapter to Chapter,
% correct umlauts, correct hyphenation)
% see also http://en.wikipedia.org/wiki/Babel-System

\usepackage{titling}
%Helps access the declared author and title in the document

\usepackage[style=numeric-comp,backend=biber,doi=false,isbn=false,maxnames=3,sorting=none]{biblatex}
\usepackage{csquotes}
\addbibresource{Referenzen.bib}
%Controls the creation of the bibliography and citation throughout the report
%At the same time accesses the Referenzen.bib file containing the sources.

\usepackage{url}
%creates nicer URLs (especially important in sources)

\usepackage{float}
%Helps arrange images, graphics and tables

%\usepackage{abstract}
%allows inserting an abstract (not relevant for internship)

\usepackage{booktabs}
%enables creation of nicer tables

\usepackage{subfigure}
%allows the user to display multiple graphics/images in one figure

\usepackage{setspace}
% Package to change line spacing
\onehalfspacing
% Set line spacing to 1.5-fold

\usepackage{pdfpages}
%enables embedding of independent PDF documents

\usepackage{siunitx}
\sisetup{locale = US}
%simplifies the required syntax for using SI units and mathematical equations with the international standard

\usepackage{tikz}
%Free arrangement of images and text in relation to each other, as well as possibilities to edit images
%Also has the ability to create graphics and diagrams in Latex

\usepackage[format=plain,labelfont=bf]{caption}
\captionsetup{position=above}
%beautifies figure and table captions. For tables, captions are displayed above the table

\usepackage{icomma}
%For correct spacing of commas in decimal numbers and in equations

\usepackage{amsmath}
%elementary package for equations. Beautifies them and facilitates the use of equations.

\usepackage{graphicx}
%more options when embedding and arranging images and graphics

\usepackage{setspace}
%For correct line spacing
%%% HJP 31/3/25
\newlength\dist
\setlength\dist{-\topskip}
\let\chapback\chapterheadstartvskip
\def\setchapspace#1{\setlength\dist{#1\dist}}
\def\lesschapspace{\renewcommand*\chapterheadstartvskip{\vspace*{\dist}}}
\def\resetchapspace{\let\chapterheadstartvskip\chapback}
%%%

\usepackage[left=3cm,right=2cm,top=3cm, bottom=3cm]{geometry}
%Controls page margins
%May need to be adjusted for double-sided layout.

\usepackage{enumitem}

\usepackage[hidelinks]{hyperref}
% Generic links

\usepackage[acronym, toc]{glossaries}

\renewcommand*{\glspostdescription}{}
%Removes numbering in the abbreviations list

\usepackage{listings}
\definecolor{codegreen}{rgb}{0,0.6,0}
\definecolor{codegray}{rgb}{0.5,0.5,0.5}
\definecolor{codepurple}{rgb}{0.58,0,0.82}
\definecolor{backcolour}{rgb}{0.95,0.95,0.92}
\lstdefinestyle{myStyle}{
    backgroundcolor=\color{backcolour},
    commentstyle=\color{codegreen},
    keywordstyle=\color{magenta},
    numberstyle=\tiny\color{codegray},
    stringstyle=\color{codepurple},
    basicstyle=\ttfamily\footnotesize,
    breakatwhitespace=false,
    breaklines=true,
    keepspaces=true,
    numbers=left,
    numbersep=5pt,
    showspaces=false,
    showstringspaces=false,
    showtabs=false,
    tabsize=2,
}

\usepackage{todonotes}
% Can be completely disabled with disable
%\usepackage[disable]{todonotes}
\setuptodonotes{fancyline}
\todototoc

\linespread{1.5}
\flushbottom

\iffalse % don’t use titlesec.sty with KOMA script classes
\usepackage{titlesec}
%% \titleformat{⟨command⟩}[⟨shape⟩]{⟨format⟩}{⟨label⟩}{⟨sep⟩}{⟨before⟩}[⟨after⟩]
\titleformat{name=\chapter}[display]
{\raggedleft\chaptertitlename\ {\textcolor{gray}{\fontsize{60}    {65}\selectfont\thechapter}}}
{20pt}{}[]
\fi

%%%%%%%%%%%%%%%%%%%%%%%%%%%%%%%%%%%%% ToC DEPTH LEVEL
\setcounter{secnumdepth}{3} % number subsubsection
\setcounter{tocdepth}{3} % list subsubsection


\addtokomafont{chapterprefix}{\raggedleft}

%%%%%%%%%%%%%%%%%%%%%%%%%%%%%%%%%%%%% PREFACE AND BODY STYLING
\def\preface{
    \pagenumbering{roman}
    \doublespacing
}

\def\body{
    %\cleardoublepage
    \lesschapspace
    \linespread{1.5}
    \pagenumbering{arabic}
    \pagestyle{headings}
}

\def\abstract{
    \begin{center}{
        \large\bfseries Abstract}
    \end{center}
    \normalsize
    \normalfont
    \linespread{1.5}
    }{\cleardoublepage}
\def\endabstract{
  \par
}

\newenvironment{acknowledgements}{
   \cleardoublepage
    \begin{center}{
        \large \bf Acknowledgements}
    \end{center}
    \normalsize
    \linespread{1.5}
    }{\cleardoublepage}
\def\endacknowledgements{
  \par
}

%%%%%%%%%%%%%%%%%%%%%%%%%%%%%%%%%%%%%%%%%%%%%%%%%%%%%%%%%
%%%%%%%%%%%%%%%%%%%%%%%%%%%%%%%%%%%%%%%%%%%%%%%%%%%%%%%%%
